% ============================================================
%  Chapter 1 -- Appendix: construction of R by Dedekind cuts
% ============================================================
\rsection{Appendix}

\begin{roadmap}[What the appendix is doing, and when to read it]
This is the promissory note from 1.19 being paid. Rudin asserted that an ordered
field with the least-upper-bound property exists; here he builds one out of
$\Q$ and nothing else.

\medskip
\raggedright
\textbf{You may skip it.} Nothing in Chapters 2--11 depends on having read it,
and Rudin says in the Preface it may be ``studied and enjoyed whenever the time
seems ripe.'' What it buys you is the knowledge that 1.19 is a theorem rather
than an article of faith. He is not alone in deferring it: Yu Pin's Tsinghua
notes call real-number theory ``very important for the completeness of the
knowledge system, but not indispensable for understanding the basic ideas of
analysis,'' and the USTC lecture notes postpone it to a third volume.

\medskip
\textbf{The idea in one sentence.} A real number is identified with the set of
all rationals below it. That set --- a \emph{cut} --- is a \emph{subset} of
$\Q$, so no new kind of object has to be conjured up; $\sqrt2$ becomes
$\{p \in \Q : p<0 \text{ or } p^2<2\}$, a perfectly ordinary collection of
rationals which exists whether or not $\sqrt 2$ does.

\medskip
\textbf{The nine steps, and which ones matter.}

\medskip
\centering
\begin{tikzpicture}[scale=0.9, transform shape, node distance=4.5mm]
\node[rmbox, text width=20mm] (a) {\textbf{1--3}\\cuts, order,\\the \lub{} property};
\node[rmbox, text width=20mm, right=of a] (b) {\textbf{4--5}\\addition};
\node[rmbox, text width=20mm, right=of b] (c) {\textbf{6--7}\\multiplication};
\node[rmbox, text width=20mm, right=of c] (d) {\textbf{8--9}\\$\Q$ sits inside};
\draw[rmarrow] (a) -- (b); \draw[rmarrow] (b) -- (c); \draw[rmarrow] (c) -- (d);
\end{tikzpicture}

\medskip
\raggedright
\textbf{Step 3 is the whole point} --- it is where the least-upper-bound
property appears, and its proof is four lines: the supremum of a set of cuts is
their \emph{union}. That is the payoff of representing numbers as sets. Steps
4--7 are bookkeeping, verifying field axioms one at a time, and Rudin openly
omits most of Step 6 as too similar to Step 4. Steps 8--9 establish the second
sentence of 1.19, that $\Q$ is a subfield.

\medskip
\textbf{If you read only two things}, read Step 1 (what a cut is, and why
property (III) is there) and Step 3.
\end{roadmap}

\noindent
Theorem 1.19 will be proved in this appendix by constructing $\R$ from $\Q$. We
shall divide the construction into several steps.

\medskip
\noindent\textbf{Step 1}\quad The members of $\R$ will be certain subsets of
$\Q$, called \emph{cuts}. A cut is, by definition, any set $\alpha \subseteq \Q$
with the following three properties.
\begin{enumerate}[label=(\Roman*), leftmargin=2.6em, itemsep=2pt, topsep=4pt]
  \item $\alpha$ is not empty, and $\alpha \ne \Q$.
  \item If $p \in \alpha$, $q \in \Q$, and $q < p$, then $q \in \alpha$.
  \item If $p \in \alpha$, then $p < r$ for some $r \in \alpha$.
\end{enumerate}
The letters $p,q,r,\dots$ will always denote rational numbers, and
$\alpha,\beta,\gamma,\dots$ will denote cuts.

Note that (III) simply says that $\alpha$ has no largest member; (II) implies
two facts which will be used freely:
\begin{quote}
If $p \in \alpha$ and $q \notin \alpha$, then $p < q$.\\
If $r \notin \alpha$ and $r < s$, then $s \notin \alpha$.
\end{quote}

\begin{unpack}[What the three properties are for]
Picture the cut $\alpha$ as ``everything strictly below some point on the
line.'' Each condition rules out a degenerate set that would spoil that picture.

\smallskip
\textbf{(I) excludes the two extremes.} $\emptyset$ and $\Q$ would correspond to
$-\infty$ and $+\infty$, which are not real numbers.

\smallskip
\textbf{(II) says $\alpha$ is downward closed} --- it is an initial segment of
$\Q$, not an arbitrary subset. Without it $\{0,5\}$ would qualify, and nothing
about order would work.

\smallskip
\textbf{(III) says $\alpha$ has no largest member}, which forces the cut to be
the set of rationals \emph{strictly} below its edge. This is the delicate one,
and its purpose is uniqueness: without it, the point $2$ would be represented by
two different cuts, $\{p<2\}$ and $\{p\le 2\}$, and the map from cuts to numbers
would be two-to-one at every rational. Exercise 1.20 asks what breaks if you
drop it --- the answer is that the additive inverse (A5) fails.

\smallskip
The two consequences of (II) that Rudin records are the workhorses of the whole
appendix. Both say the same thing --- a cut has no gaps and no stragglers ---
and he uses them dozens of times without further comment.
\end{unpack}

\begin{center}
\begin{tikzpicture}[x=1mm, y=1mm]
  % a cut
  \draw[axis] (-1,10) -- (92,10);
  \draw[seg] (2,10) -- (52,10);
  \node[ptopen] at (52,10) {};
  \node[lbl, anchor=north] at (26,8) {$\alpha$ --- everything in $\Q$ below the edge};
  \node[lbl, anchor=south] at (74,11) {$\Q\setminus\alpha$};
  \node[lbl, anchor=south] at (52,13.5) {no largest member: (III)};
  \draw[guide] (52,11.5) -- (52,13);
  % a second, larger cut
  \draw[axis] (-1,-8) -- (92,-8);
  \draw[seg] (2,-8) -- (70,-8);
  \node[ptopen] at (70,-8) {};
  \node[lbl, anchor=north] at (30,-10) {$\beta \supset \alpha$, so $\alpha<\beta$ (Step 2)};
  \draw[guide] (52,9) -- (52,-8);
\end{tikzpicture}
\end{center}
\figcap{A cut is a leftward ray of $\Q$ with no endpoint. Order is containment;
one cut is smaller when it is a proper subset.}

\medskip
\noindent\textbf{Step 2}\quad Define ``$\alpha < \beta$'' to mean: $\alpha$ is a
proper subset of $\beta$.

Let us check that this meets the requirements of Definition 1.5.

If $\alpha < \beta$ and $\beta < \gamma$ it is clear that $\alpha < \gamma$. (A
proper subset of a proper subset is a proper subset.) It is also clear that at
most one of the three relations
\[ \alpha < \beta, \qquad \alpha = \beta, \qquad \beta < \alpha \]
can hold for any pair $\alpha, \beta$. To show that at least one holds, assume
that the first two fail. Then $\alpha$ is not a subset of $\beta$. Hence there
is a $p \in \alpha$ with $p \notin \beta$. If $q \in \beta$, it follows that
$q < p$ (since $p \notin \beta$), hence $q \in \alpha$, by (II). Thus
$\beta \subset \alpha$. Since $\beta \ne \alpha$, we conclude
$\beta < \alpha$.\kn{Trichotomy is the only part needing work, because
containment of arbitrary sets is a \emph{partial} order --- two sets can be
incomparable. What rules that out here is (II): downward-closed subsets of a
totally ordered set are automatically nested. This is the first place property
(II) earns its keep.}

Thus $\R$ is now an ordered set.

\medskip
\noindent\textbf{Step 3}\quad \emph{The ordered set $\R$ has the least upper
bound property.}

To prove this, let $A$ be a nonempty subset of $\R$, and assume that
$\beta \in \R$ is an upper bound of $A$. Define $\gamma$ to be the union of all
$\alpha \in A$. In other words, $p \in \gamma$ if and only if $p \in \alpha$ for
some $\alpha \in A$. We shall prove that $\gamma \in \R$ and that
$\gamma = \sup A$.

Since $A$ is not empty, there exists an $\alpha_0 \in A$. This $\alpha_0$ is not
empty. Since $\alpha_0 \subset \gamma$, $\gamma$ is not empty. Next,
$\gamma \subset \beta$ (since $\alpha \subset \beta$ for every $\alpha \in A$),
and therefore $\gamma \ne \Q$. Thus $\gamma$ satisfies property (I). To prove
(II) and (III), pick $p \in \gamma$. Then $p \in \alpha_1$ for some
$\alpha_1 \in A$. If $q < p$, then $q \in \alpha_1$, hence $q \in \gamma$; this
proves (II). If $r \in \alpha_1$ is so chosen that $r > p$, we see that
$r \in \gamma$ (since $\alpha_1 \subset \gamma$), and therefore $\gamma$
satisfies (III). Thus $\gamma \in \R$.

It is clear that $\alpha \le \gamma$ for every $\alpha \in A$.

Suppose $\delta < \gamma$. Then there is an $s \in \gamma$ such that
$s \notin \delta$. Since $s \in \gamma$, $s \in \alpha$ for some $\alpha \in A$.
Hence $\delta < \alpha$ and $\delta$ is not an upper bound of $A$.

This gives the desired result: $\gamma = \sup A$.

\begin{deeper}[Why both hypotheses of 1.10 appear, and where]
Watch where each is consumed. \emph{Nonempty} gives an $\alpha_0 \in A$, which
is what makes $\gamma$ nonempty --- half of property (I). \emph{Bounded above}
gives a $\beta$ containing every $\alpha$, hence $\gamma \subset \beta \ne \Q$
--- the other half.

Without either, the union might not be a cut: an empty union is $\emptyset$, and
an unbounded union is all of $\Q$. Both are excluded by (I) precisely because
they are the cuts that would represent $-\infty$ and $+\infty$. This is the same
observation as the warning at 1.10, now visible from the inside: the two
hypotheses are not conveniences, they are what keeps the answer inside the
system.

The last two paragraphs are the two clauses of Definition 1.8: $\gamma$ is an
upper bound, and nothing smaller is. Note how cheap the second is --- if
$\delta$ is properly contained in the union, it must miss an element that some
$\alpha$ contains, so $\alpha$ already exceeds it.
\end{deeper}

\medskip
\noindent\textbf{Step 4}\quad If $\alpha \in \R$ and $\beta \in \R$, we define
$\alpha + \beta$ to be the set of all sums $r+s$, where $r \in \alpha$ and
$s \in \beta$.

We define $0^*$ to be the set of all negative rational numbers. It is clear that
$0^*$ is a cut. We verify that \emph{the axioms for addition} (see Definition
1.12) \emph{hold in $\R$, with $0^*$ playing the role of $0$.}

\begin{enumerate}[label=(A\arabic*), leftmargin=3.1em, itemsep=6pt, topsep=5pt]
  \item We have to show that $\alpha+\beta$ is a cut. It is clear that
        $\alpha+\beta$ is a nonempty subset of $\Q$. Take $r' \notin \alpha$,
        $s' \notin \beta$. Then $r'+s' > r+s$ for all choices of $r \in \alpha$,
        $s \in \beta$. Thus $r'+s' \notin \alpha+\beta$. It follows that
        $\alpha+\beta$ has property (I).

        Pick $p \in \alpha+\beta$. Then $p = r+s$, with $r \in \alpha$,
        $s \in \beta$. If $q < p$, then $q - s < r$, so $q - s \in \alpha$, and
        $q = (q-s)+s \in \alpha+\beta$. Thus (II) holds. Choose $t \in \alpha$
        so that $t > r$. Then $p < t+s$ and $t+s \in \alpha+\beta$. Thus (III)
        holds.
  \item $\alpha+\beta$ is the set of all $r+s$, with $r \in \alpha$,
        $s \in \beta$. By the same definition, $\beta+\alpha$ is the set of all
        $s+r$. Since $r+s = s+r$ for all $r \in \Q$, $s \in \Q$, we have
        $\alpha+\beta = \beta+\alpha$.
  \item As above, this follows from the associative law in $\Q$.
  \item If $r \in \alpha$ and $s \in 0^*$, then $r+s < r$, hence
        $r+s \in \alpha$. Thus $\alpha+0^* \subset \alpha$. To obtain the
        opposite inclusion, pick $p \in \alpha$, and pick $r \in \alpha$,
        $r > p$. Then $p - r \in 0^*$, and
        $p = r+(p-r) \in \alpha+0^*$. Thus $\alpha \subset \alpha+0^*$. We
        conclude that $\alpha+0^* = \alpha$.\kn{The reverse inclusion is where
        property (III) is used, and it is used essentially: we need some
        $r \in \alpha$ strictly \emph{above} $p$ so that $p-r$ is negative. If
        $\alpha$ had a largest member $p$, no such $r$ would exist and
        $\alpha + 0^*$ would be missing $p$.}
  \item Fix $\alpha \in \R$. Let $\beta$ be the set of all $p$ with the
        following property:
        \begin{quote}
        There exists $r > 0$ such that $-p-r \notin \alpha$.
        \end{quote}
        In other words, some rational number smaller than $-p$ fails to be in
        $\alpha$.

        We show that $\beta \in \R$ and that $\alpha+\beta = 0^*$.

        If $s \notin \alpha$ and $p = -s-1$, then $-p-1 \notin \alpha$, hence
        $p \in \beta$. So $\beta$ is not empty. If $q \in \alpha$, then
        $-q \notin \beta$. So $\beta \ne \Q$. Hence $\beta$ satisfies (I).

        Pick $p \in \beta$ and $r > 0$ so that $-p-r \notin \alpha$. If
        $q < p$, then $-q-r > -p-r$, hence $-q-r \notin \alpha$. Thus
        $q \in \beta$ and (II) holds. Put $t = p+(r/2)$. Then $t > p$ and
        $-t-(r/2) = -p-r \notin \alpha$, so that $t \in \beta$. Hence $\beta$
        satisfies (III).

        We have proved that $\beta \in \R$.

        If $r \in \alpha$ and $s \in \beta$, then $-s \notin \alpha$, hence
        $r < -s$, $r+s < 0$. Thus $\alpha+\beta \subset 0^*$.

        To prove the opposite inclusion, pick $v \in 0^*$, put $w = -v/2$. Then
        $w > 0$ and there is an integer $n$ such that $nw \in \alpha$ but
        $(n+1)w \notin \alpha$. (Note that this depends on the fact that $\Q$
        has the archimedean property!) Put $p = -(n+2)w$. Then $p \in \beta$,
        since $-p-w \notin \alpha$, and
        \[ v = nw + p \in \alpha+\beta. \]
        Thus $0^* \subset \alpha+\beta$.

        We conclude that $\alpha+\beta = 0^*$. This $\beta$ will of course be
        denoted by $-\alpha$.
\end{enumerate}

\begin{unpack}[Why $-\alpha$ is defined so awkwardly]
The natural guess is that $-\alpha$ should be $\{-q : q \notin \alpha\}$ ---
reflect the complement. That is almost right and fails on exactly one point.

If $\alpha$ is the cut of a rational $r$, then $\Q \setminus \alpha$ has a
smallest element, namely $r$ itself. Reflecting gives a set with a
\emph{largest} element $-r$, violating property (III). So the naive definition
does not always produce a cut.

Rudin's ``there exists $r>0$ with $-p-r \notin \alpha$'' repairs this by
demanding strict room to spare: $p$ qualifies only if something \emph{strictly
below} $-p$ already lies outside $\alpha$. That extra $r$ shaves off the single
offending endpoint and nothing else. You can see the repair working in the
verification of (III), where $t = p + r/2$ is produced by spending half the
slack.

The archimedean appeal at the end is the other subtlety, and Rudin flags it
himself with an exclamation mark. To hit an arbitrary $v \in 0^*$ we need
integer multiples of $w$ to eventually escape $\alpha$; that is exactly 1.20(a),
holding in $\Q$. Note the logical order --- the archimedean property of $\Q$ is
elementary and assumed, and it is being used to build $\R$, not derived from it.
\end{unpack}

\medskip
\noindent\textbf{Step 5}\quad Having proved that the addition defined in Step 4
satisfies axioms (A) of Definition 1.12, it follows that Proposition 1.14 is
valid in $\R$, and we can prove one of the requirements in Definition 1.17:
\begin{quote}
If $\alpha,\beta,\gamma \in \R$ and $\beta < \gamma$, then
$\alpha+\beta < \alpha+\gamma$.
\end{quote}
Indeed, it is obvious from the definition of $+$ in $\R$ that
$\alpha+\beta \subset \alpha+\gamma$; if we had $\alpha+\beta = \alpha+\gamma$,
the cancellation law (Proposition 1.14) would imply $\beta = \gamma$.\kd{Exactly
the dividend promised at 1.13(c). Proposition 1.14 was proved once, for an
arbitrary set satisfying axioms (A); having checked those axioms in Step 4,
Rudin may now invoke cancellation for cuts without reproving anything. The
one-time purchase is being collected.}

It also follows that $\alpha > 0^*$ if and only if $-\alpha < 0^*$.

\medskip
\noindent\textbf{Step 6}\quad Multiplication is a little more bothersome than
addition in the present context, since products of negative rationals are
positive. For this reason we confine ourselves first to $\R^+$, the set of all
$\alpha \in \R$ with $\alpha > 0^*$.

If $\alpha \in \R^+$ and $\beta \in \R^+$, we define $\alpha\beta$ to be the set
of all $p$ such that $p \le rs$ for some choice of $r \in \alpha$,
$s \in \beta$, $r>0$, $s>0$.\kw{The definition is \emph{not} ``the set of all
products $rs$'' --- that would fail property (II), since the products of
positive rationals omit everything negative. The clause $p \le rs$ patches this
by throwing in everything below. Contrast Step 4, where the sums $r+s$ already
sweep out a downward-closed set on their own.}

We define $1^*$ to be the set of all $q < 1$.

Then \emph{the axioms} (M) \emph{and} (D) \emph{of Definition 1.12 hold, with
$\R^+$ in place of $F$, and with $1^*$ in the role of $1$.}

The proofs are so similar to the ones given in detail in Step 4 that we omit
them.

Note, in particular, that the second requirement of Definition 1.17 holds: If
$\alpha > 0^*$ and $\beta > 0^*$, then $\alpha\beta > 0^*$.

\medskip
\noindent\textbf{Step 7}\quad We complete the definition of multiplication by
setting $\alpha 0^* = 0^*\alpha = 0^*$, and by setting
\[
  \alpha\beta = \begin{cases}
    (-\alpha)(-\beta) & \text{if } \alpha<0^*,\ \beta<0^*,\\
    -[(-\alpha)\beta] & \text{if } \alpha<0^*,\ \beta>0^*,\\
    -[\alpha(-\beta)] & \text{if } \alpha>0^*,\ \beta<0^*.
  \end{cases}
\]
The products on the right were defined in Step 6.

Having proved (in Step 6) that the axioms (M) hold in $\R^+$, it is now
perfectly simple to prove them in $\R$, by repeated application of the identity
$\gamma = -(-\gamma)$, which is part of Proposition 1.14. (See Step 5.)

The proof of the distributive law
\[ \alpha(\beta+\gamma) = \alpha\beta + \alpha\gamma \]
breaks into cases. For instance, suppose $\alpha > 0^*$, $\beta < 0^*$,
$\beta+\gamma > 0^*$. Then $\gamma = (\beta+\gamma)+(-\beta)$, and (since we
already know that the distributive law holds in $\R^+$)
\[ \alpha\gamma = \alpha(\beta+\gamma) + \alpha\cdot(-\beta). \]
But $\alpha\cdot(-\beta) = -(\alpha\beta)$. Thus
\[ \alpha\beta + \alpha\gamma = \alpha(\beta+\gamma). \]
The other cases are handled in the same way.

We have now completed the proof that \emph{$\R$ is an ordered field with the
least upper bound property.}

\begin{deeper}[Why multiplication needs cases and addition does not]
Addition is monotone in each argument on all of $\Q$: if $r$ grows, so does
$r+s$. Multiplication is not --- multiplying by a negative reverses order. Since
a cut is defined by \emph{downward} closure, the union of products
$\{rs\}$ points the wrong way as soon as a negative factor appears, and the
Step 6 definition simply does not describe a cut.

Rudin's escape is to define multiplication only where it is well behaved, on
$\R^+$, and then extend by the sign rules --- which is the same device used at
1.16(c),(d) to prove $(-x)(-y) = xy$ from the field axioms. The cases in Step 7
are those rules read as definitions rather than theorems.

This asymmetry is why the appendix is longer than it looks. It is also why the
distributive law, which mixes the two operations, must be checked case by case:
$\beta$ and $\gamma$ can have opposite signs while $\beta+\gamma$ has either.
\end{deeper}

\medskip
\noindent\textbf{Step 8}\quad We associate with each $r \in \Q$ the set $r^*$
which consists of all $p \in \Q$ such that $p < r$. It is clear that each $r^*$
is a cut; that is, $r^* \in \R$. These cuts satisfy the following relations:
\begin{enumerate}[label=(\alph*), leftmargin=2.2em, itemsep=1pt, topsep=3pt]
  \item $r^* + s^* = (r+s)^*$,
  \item $r^*s^* = (rs)^*$,
  \item $r^* < s^*$ if and only if $r < s$.
\end{enumerate}

To prove (a), choose $p \in r^*+s^*$. Then $p = u+v$, where $u<r$, $v<s$. Hence
$p < r+s$, which says that $p \in (r+s)^*$.

Conversely, suppose $p \in (r+s)^*$. Then $p < r+s$. Choose $t$ so that
$2t = r+s-p$, put
\[ r' = r-t, \qquad s' = s-t. \]
Then $r' \in r^*$, $s' \in s^*$, and $p = r'+s'$, so that $p \in
r^*+s^*$.\kn{The trick is splitting the shortfall evenly. We know $p$ falls
short of $r+s$ by $2t>0$; docking $t$ from each of $r$ and $s$ keeps both
strictly below their targets while making the sum land exactly on $p$. Halving
the slack is the same device as $t = p + r/2$ in Step 4(A5).} This proves (a).
The proof of (b) is similar.

If $r < s$ then $r \in s^*$, but $r \notin r^*$; hence $r^* < s^*$.

If $r^* < s^*$, then there is a $p \in s^*$ such that $p \notin r^*$. Hence
$r \le p < s$, so that $r < s$. This proves (c).

\medskip
\noindent\textbf{Step 9}\quad We saw in Step 8 that the replacement of the
rational numbers $r$ by the corresponding ``rational cuts'' $r^* \in \R$
preserves sums, products, and order. This fact may be expressed by saying that
the ordered field $\Q$ is \emph{isomorphic} to the ordered field $\Q^*$ whose
elements are the rational cuts. Of course, $r^*$ is by no means the same as $r$,
but the properties we are concerned with (arithmetic and order) are the same in
the two fields.

It is this identification of $\Q$ with $\Q^*$ which allows us to regard $\Q$ as
a subfield of $\R$.\kd{This settles the second sentence of Theorem 1.19, and it
is worth being clear about what is and is not claimed. The rational $2$ and the
cut $2^* = \{p \in \Q : p<2\}$ are genuinely different objects --- one a number,
the other an infinite set. What Step 8 establishes is that no
arithmetic-or-order question can tell them apart, and we therefore agree to
write $\Q \subset \R$. The same courtesy was extended to $\R \subset \C$ after
1.26.}

The second part of Theorem 1.19 is to be understood in terms of this
identification. Note that the same phenomenon occurs when the real numbers are
regarded as a subfield of the complex field, and it also occurs at a much more
elementary level, when the integers are identified with a certain subset of
$\Q$.

It is a fact, which we will not prove here, that any two ordered fields with the
least-upper-bound property are isomorphic. The first part of Theorem 1.19
therefore \emph{characterizes} the real field $\R$ completely.\kd{The strongest
sentence in the appendix, and it is dropped without proof. It says that
Dedekind's construction was not one choice among many with different outcomes:
\emph{any} route to a complete ordered field --- cuts, Cauchy sequences,
infinite decimals --- lands on the same object up to isomorphism. That is why
nothing in the rest of the book ever refers back to cuts. $\R$ is fully
determined by its properties, so the properties are all one needs.}

The books by Landau and Thurston cited in the Bibliography are entirely devoted
to number systems. Chapter 1 of Knopp's book contains a more leisurely
description of how $\R$ can be obtained from $\Q$. Another construction, in
which each real number is defined to be an equivalence class of Cauchy sequences
of rational numbers (see Chap.\ 3), is carried out in Sec.\ 5 of the book by
Hewitt and Stromberg.

The cuts in $\Q$ which we used here were invented by Dedekind. The construction
of $\R$ from $\Q$ by means of Cauchy sequences is due to Cantor. Both Cantor and
Dedekind published their constructions in 1872.\kd{Cantor's route is the one
most modern texts take --- Tao's \emph{Analysis I} builds $\R$ as equivalence
classes of Cauchy sequences of rationals (his Chapter 5) and explicitly declines
Dedekind cuts. It generalises: the same construction completes any metric space
(Rudin does exactly this in Exercise 3.24), whereas cuts need an order and so
stop at $\R$. Cuts are shorter here, which is why Rudin keeps them. Both books
warn off the third route, infinite decimals: it looks familiar but has to
legislate $0.999\dots = 1$, and is the most delicate of the three.}

\begin{recap}[What the appendix established]
$\R$ exists. Concretely, it is the set of downward-closed, unbounded-below,
maximum-free subsets of $\Q$, ordered by inclusion, with addition given by
elementwise sums and multiplication defined first for positives and then by sign
rules. Under those definitions it is an ordered field, its least upper bounds are
unions, and it contains an isomorphic copy of $\Q$.

\smallskip
Theorem 1.19 is therefore a theorem. And since complete ordered fields are
unique up to isomorphism, this particular construction may now be forgotten ---
which is precisely what the rest of the book does.
\end{recap}
